% =============================================================================
%  Sección 2: Patrimonio
% =============================================================================

\section{Fase 2: Mapeo del Patrimonio (Tangible e Intangible)}
\label{sec:mapeo-patrimonial}

El mapeo patrimonial de San Pedro Atocpan trasciende la catalogación monumental; 
se presenta como una \termino{Puesta en Valor} que articula la arquitectura religiosa con los 
procesos de resiliencia social derivados de la tragedia hídrica de 1935.

\subsection{Inventario Arquitectónico (Patrimonio Tangible)}
\label{ssec:inventario-arquitectonico}

Siguiendo los lineamientos de edificios de primer orden, se identifican tres inmuebles 
clave dentro de los polígonos definidos y el pueblo de San Pedro Atocpan:

\begin{itemize}
    \item \textbf{Parroquia de San Pedro Apóstol:} Centro neurálgico fundacional cuyo nombre 
            oficial data de 1532. Representa el eje de la traza urbana tradicional.
    \item \textbf{Capilla de San Martín Caballero (Yencuictlalpan):} Localizada en el 
            \textit{Polígono 2}, es el recinto histórico que resguardaba 
            la figura del Señor de las Misericordias durante la tromba de 1935. 
            Su atrio sirvió como refugio ante la ola de lodo y agua que alcanzó 
            los 2 metros de altura  \cite{atocpan2026relato}.
    \item \textbf{Santuario del Señor de las Misericordias:} Ubicado en el 
            \textit{Polígono 3}, es una obra de arquitectura religiosa moderna concluida 
            a finales de la década de 1960. Su construcción es un testimonio de la soberanía comunal, 
            al ser financiada íntegramente por las aportaciones de los habitantes \cite{santuarioSFtesoro}.
\end{itemize}

\subsection{Tejido Social e Identidad (Patrimonio Intangible)}
\label{ssec:tejido-social}

La identidad de Atocpan está anclada en una memoria colectiva que amalgama la tradición 
étnica con la superación de desastres naturales.

\subsubsection{Memoria Oral y Resiliencia}
\label{subsubsec:memoria-oral}

A falta de registros oficiales extensos, la memoria de la "Gran Tromba" sobrevive a través 
de testimonios como el del Sr. Conrado Amaya Cabello (\textit{Tatitla Tontle}), 
quien conserva pasajes del desastre y la posterior visita de auxilio del general 
Lázaro Cárdenas. Esta tradición oral es un pilar de la cohesión social 
que define el sentido de pertenencia en el barrio 
\cite{atocpan2026relato,excelsior1935cuatrocientas1,diccionario_nahuatl2018}.

\subsubsection{Cultura Náhuatl y Toponimia}
\label{subsubsec:cultura-nahuatl}

La persistencia de la lengua náhuatl se manifiesta en la vida cotidiana y en 
situaciones de crisis, como los gritos de auxilio \textit{“Xinechmaka momaj”} 
(Dame tu mano) documentados durante la inundación. 
La propia toponimia de \textit{Atocpan} \termino{Sobre el lodo} es 
una herencia lingüística que advierte sobre la naturaleza geográfica del sitio.

\subsubsection{Dinámica Socioeconómica: La Capital del Mole}
\label{subsubsec:dinamica-mole}
El pueblo ha transformado su vocación agrícola tradicional en una potencia 
comercial gastronómica. La producción de mole y la \textit{Feria Internacional del Mole} 
(celebrada desde 1977 en el sector Momoxco) representan 
el motor económico actual y una de las manifestaciones culturales más 
potentes de la región.

% Bloque sugerido para imagen de identidad local
\begin{figure}[H]
    \centering
    %\includegraphics[width=0.7\textwidth]{img/feria_mole.png}
    \caption{Inauguración de la Feria Nacional del Mole, símbolo del patrimonio 
    económico e intangible de San Pedro Atocpan. Fuente: \cite{cabrera2025tromba}.}
    \label{fig:patrimonio_social}
\end{figure}